\Needspace{2in}
\begin{prob}
    Solve the below recurrence, stating the solution in asymptotic notation. You do not
    need to show your work (and your work won't be graded for partial credit, so be
    careful!).
    \[
        T(n) =
        \begin{cases}
            T(n/3) + \Theta(n) & n > 1 \\
            \Theta(1) & n = 1
        \end{cases}
    \]

    \inlineresponsebox[.75in]{$\Theta(n)$}

    \begin{soln}
        $\Theta(n)$.

        Unrolling several times, we find:

        \begin{align*}
            T(n)
            &=
                T(n/3) + n
                \\
            &=
                T(n/9) + \frac{n}{3} + n
                \\
            &=
                T(n/27) + \frac{n}{9} + \frac{n}{3} + n
                \\
        \end{align*}

        On the $k$th unroll, we'll get:
        \[
            T(n) = T(n/3^k) + \left[
                n + \frac{n}{3} + \frac{n}{9} + \ldots + \frac{n}{3^{k-1}}
            \right]
        \]

        It will take $\log_3 n$ unrolls to each the base case. Substituting this for
        $k$, we get:
        \begin{align*}
            T(n)
            &= T(n/3^{\log_3 n}) + \left[
                n + \frac{n}{3} + \frac{n}{9} + \ldots + \frac{n}{3^{(\log_3 n)-1}}
            \right]
            \\
            &= T(1) + n
            \underbrace{\left[
                1 + \frac{1}{3} + \frac{1}{9} + \ldots + \frac{1}{3^{(\log_3 n)-1}}
            \right]}_{\leq 2}
        \intertext{Here we've recognized the summation in square braces as a geometric
            sum of powers of $1/3$, which converges to 2 when there are infinitely-many
            terms. We've seen this sum in solving a couple of homework problems.}
            & = T(1) + c n\\
            & = \Theta(1) + c n\\
            &= \Theta(n)
        \end{align*}

    \end{soln}
\end{prob}
